\section{总体分析}
\subsection{数据特征分析}\label{sec:data}

\subsubsection{空间与地形}

16 个任务节点（1 个调度中心 O01 与 15 个服务区）的坐标为 WGS84 经纬度，节点地面海拔介于 \SIrange{127.7}{444.5}{m}。由于\textbf{经纬度直接计算欧氏距离会产生显著误差}，本文统一先投影到米制平面坐标系再计算水平距离，并断言投影后单位为米。

区域 30\,m DEM 覆盖经度 $109.03^\circ\sim109.45^\circ$、纬度 $22.86^\circ\sim23.22^\circ$。需要特别说明的是，附件 DEM 为\textbf{数字表面模型（DSM，含植被与建筑）}，本文按原样使用、不做“去建筑”处理，因此\textbf{净空与遮挡判定结果偏保守}，对安全性有利。

\subsubsection{物资与时限}

80 个货箱共 \SI{758}{kg}、\SI{2.011}{m^3}，按物资类型分为 4 类。两类时限必须严格区分：\textbf{首批截止时间}仅约束标记为“首批保障”的 30 个货箱（每个服务区恰好 2 箱：医疗 1 $+$ 饮用水 1），是\textbf{硬约束}；\textbf{期望送达时间}覆盖全部 80 箱，其中医疗物资的期望送达亦按硬约束处理，其余物资的期望时间用于衡量配送及时性。

\subsubsection{装备与链路}

三种运输机型的核心参数如表~\ref{tab:uav} 所示。附件\textbf{只给出空载/满载标准航程与电池可用能量，未给出运输机巡航功率}（仅中继机给出功率），这一点直接决定了 \S\ref{sec:energy} 的能耗口径必须“由航程反推”。

\begin{table}[htbp]
	\caption{三种运输机型主要参数（附件直接核对）}
	\label{tab:uav}
	\centering
	\zihao{-5}
	\begin{tabular}{lccc}
		\toprule
		参数 & A 型 & B 型 & C 型 \\
		\midrule
		最大载货质量 $\Qg$ / kg            & 25   & 30   & 80   \\
		可用装载体积 $V_{g}$ / $\mathrm{m^3}$ & 0.060 & 0.073 & 0.250 \\
		空载标准航程 $L_{g0}$ / km          & 25   & 28   & 26   \\
		满载标准航程 $L_{gF}$ / km          & 20   & 16   & 12   \\
		电池可用能量 $\Euse$ / kWh          & 4.5  & 4.0  & 8.0  \\
		巡航速度 $v_{g}^{c}$ / $\mathrm{m\,s^{-1}}$   & 12   & 15   & 15   \\
		爬升速度 $v_{g}^{\uparrow}$ / $\mathrm{m\,s^{-1}}$ & 3.0 & 3.0 & 2.5 \\
		下降速度 $v_{g}^{\downarrow}$ / $\mathrm{m\,s^{-1}}$ & 2.5 & 2.5 & 2.0 \\
		爬升能耗效率 $\eta_{up}$            & 0.72 & 0.72 & 0.72 \\
		返航安全余量 $\rhog$                & 20\% & 20\% & 20\% \\
		共享电池组数                        & 6    & 4    & 4    \\
		等效完全充电时间 $\Tfull$ / s       & 1800 & 2400 & 3000 \\
		\bottomrule
	\end{tabular}
\end{table}

实体机队为 A 型 4 架（U01--U04）、B 型 2 架（U05--U06）、C 型 2 架（U07--U08），共 8 架；中继无人机 2 架（R01--R02）配 6 组能源组件。中继机关键参数：起飞总质量 \SI{23.5}{kg}、巡航功率 \SI{1.15}{kW}、悬停功率 \SI{1.05}{kW}、通信附加功率 \SI{0.05}{kW}、可用能量 \SI{3.2}{kWh}、悬停离地高度上限 \SI{300}{m}、准备/建链/周转时间分别为 \SI{180}{/}\SI{30}{/}\SI{300}{s}。

链路预算参数：载波 \SI{2400}{MHz}、系统损耗 \SI{3}{dB}、遮挡附加损耗 \SI{10}{dB}、接收灵敏度 \SI{-98}{dBm}、衰落裕量 \SI{8}{dB}；网关天线离地 \SI{20}{m}。由此算得三条链路的\textbf{双向门限}（两方向取较小值）：直连 \SI{122}{dB}、中继接入 \SI{116}{dB}、中继回传 \SI{126}{dB}。

\subsection{总体技术路线}

围绕四个层层递进的问题，本文确立“\textbf{统一口径 $\to$ 逐问建模 $\to$ 独立校验}”的技术路线：先构建四问共用的公共物理与通信模型（\S\ref{sec:common}），杜绝口径分叉；再按问题一至问题四依次建模求解（\S\ref{sec:q1}--\S\ref{sec:q4}），每一问都显式给出\textbf{模型名称、决策变量、约束、求解方法与结果}；最后由独立校验器（\S\ref{sec:verify}）从附件与 DEM 重算全部物理量并与方案上报值比对。

四问之间存在严格的\textbf{数据继承关系}：问题二继承问题一的载荷口径与组批思想，但允许跨服务区串飞；问题三继承问题二的\textbf{完整运输方案}（运输架次数与能耗必须逐位一致）；问题四继承问题三的\textbf{架次划分}（“同架次的服务区必须同组”）。这种链式继承使全文结果可以互相印证，也意味着\textbf{上游任何一处修正都必须向下游传导重算}——本文在 \S\ref{sec:changes} 记录了两次这样的传导。
